\section{The Gate-Aware Adversary}
\label{sec:gateaware}

Every adversary so far has been oblivious to the defence. This one is not. It is
given the published calibration constant---the honest-coincidence ceiling for the
answer space---and it is asked to do the obvious thing: coordinate as much as it
can without crossing it.

\subsection{Construction}

A bloc coordinating with probability $p$ and otherwise picking a wrong answer
uniformly realises $q(p) = p^2 + (1-p)^2/(C-1)$, so sweeping $p$ sweeps $q$ and
the adversary picks the $p$ whose $q$ does most damage while staying below the
ceiling. The map is not monotone --- $q$ falls as $p$ rises from zero, reaches an
interior minimum, then rises --- which is why the adversary has a genuine optimum
rather than a corner solution, and why Figure~\ref{fig:evasion} plots points and
not a curve.

\subsection{The anchor is $f = \numEvasionAnchorF$}

The anchor is the exact tie, and the choice has to be defended because the more
dramatic numbers live elsewhere. At $f=0.7$ the Byzantine bloc is an outright
majority and \emph{every} coherence level, independent lying included, drives
accuracy toward zero; a collapse there cannot be attributed to the gate, and an
evasion effect is not separable from numerical dominance. An earlier version of
this work reported the $f=0.7$ attacker gain; it is withdrawn. At
$f=\numEvasionAnchorF$ the outcome turns on whether the gate fires --- and, at
that point only, on the normalisation of the adversary's own vote
(Section~\ref{sec:evaluation}), so the sweep runs in the matched-parity condition
and the confound is controlled rather than assumed away.

\subsection{Result}

The gate-aware adversary defeats AIP. At its optimum, $\hat q =
\numEvasionQMedqa$, AIP-gated falls to $\numEvasionAccMedqa$ on MedQA against
$\numEvasionCoherentMedqa$ under a fully coherent adversary---an attacker gain of
$\numEvasionGainMedqa$, eight times the floor. On MMLU the same optimum yields
$\numEvasionAccMmlu$ against $\numEvasionCoherentMmlu$, a gain of
$\numEvasionGainMmlu$. \textbf{[G1]}

The number that makes this more than an engineering setback is the optimum's
location. The chance coherence for a four-option space is $1/(C-1) =
\numEvasionChance$. The attacker's optimum sits at $\numEvasionQMedqa$:
\textbf{below chance}. At its most damaging the adversary is not merely quieter
than the calibrated ceiling---it is \emph{less} coherent than honest agents
failing independently. Raising, lowering, smoothing or randomising a threshold
cannot help, because there is no threshold on this axis that separates the
attacker from honest disagreement. Section~\ref{sec:limits} makes that argument
precisely.

MATH-500 behaves differently and the difference is instructive. There the
attacker's global optimum is a \emph{fully coherent} channel, not an evading one:
the open-text ceiling of the answer space is high enough relative to achievable
coherence that the gate is near-inert, and the adversary has nothing to evade.
Its best below-ceiling attack yields $\numEvasionAccMathfive$, which is
\emph{worse} for the attacker than coherence---a gain of
$\numEvasionGainMathfive$. A low ceiling is not a structural defence.
\textbf{[L1]}

\subsection{Pre-registered mitigations}

Three mitigations were specified, with success criteria committed to version
control before any of them ran.

\textbf{M1, soft gating.} Replace the hard binomial threshold with a logistic
function of the coherence ratio, so that a channel just under the ceiling is
partially inverted rather than fully trusted. Within the evasion band---the cells
below the ceiling, which is where the residual attack lives---mean accuracy rises
from $\numMitigHardBandMean$ to $\numMitigSoftBandMean$ and the worst case from
$\numMitigHardBandWorst$ to $\numMitigSoftBandWorst$. This is a real improvement
of about one and a half times the floor.

Over the \emph{whole} attack space at the anchor, however, the improvement is
$\numMitigHardMean$ to $\numMitigSoftMean$---exactly at the floor. Only the
band-restricted statement is defensible, and M1 fails its pre-registered success
criterion regardless: it does not bring the worst case anywhere near the
committed target. \textbf{[D1]}

\textbf{M2, threshold randomisation. Refuted, and undetectably so.} Randomising
the threshold within a relative window was intended to deny the adversary a fixed
target. It reads $\numMitigRandMean$ against the hard gate's
$\numMitigHardMean$ over the whole space and $\numMitigRandBandMean$ against
$\numMitigHardBandMean$ within the band. Both differences are far below the
floor. Randomising a statistic that carries no signal buys no discrimination---in
either direction. An earlier reading of this experiment called M2 ``actively
harmful'' on the strength of a difference that was itself below the floor; that
reading was wrong in the same way, and is withdrawn. \textbf{[D2]}

\textbf{M3, both together.} Band mean $\numMitigBothBandMean$, worst case
$\numMitigBothBandWorst$. Not separable from M1 alone.

\textbf{Both M1 and M2 were initially inert}, returning results identical to the
unmitigated baseline, and those identical results were briefly read as a finding.
They were wiring bugs, caught by probe tests that assert variant activity before
any verdict is read, and two published verdicts were retracted;
Section~\ref{sec:limitations} gives the detail and why it belongs in the paper.
\textbf{[R1]}
